\section{Fundamentos y Aprendizaje Autónomo}
\label{sec:fundamentos}

\subsection{Conceptos y Algoritmos Requeridos}
El desarrollo del proyecto descansa sobre tres pilares conceptuales de las ciencias de la computación: la teoría de grafos viales dirigidos, el diseño de colas de prioridad con clave disminuible y la relajación voraz (\textit{greedy}) de caminos mínimos.

\subsubsection{Grafos Viales Urbanos y Propiedad de Red Plana}
Una red vial urbana se modela formalmente como un grafo dirigido ponderado $G = (V, E, W)$, donde $V$ representa el conjunto de intersecciones viales o discontinuidades geométricas, $E \subseteq V \times V$ representa los tramos viales orientados según el sentido de circulación reglamentario, y $W: E \to \mathbb{R}^+$ asigna un costo no negativo a cada tramo. Las redes viales urbanas presentan propiedades geométricas cuasi-planas con una densidad extremadamente baja:
\begin{equation}
    |E| \approx k \cdot |V|, \quad \text{donde } 2.0 \le k \le 3.5
\end{equation}
En el recorte vial de OpenStreetMap de Cusco procesado en este proyecto ($|V|=891, |E|=2139$), la razón es $|E|/|V| = 2.40$. Esta dispersión estructural condiciona de forma determinante el comportamiento asintótico y práctico de los algoritmos de optimización de rutas.

\subsubsection{Algoritmo de Dijkstra e Invariante de Subestructura Óptima}
El algoritmo de Dijkstra \parencite{dijkstra1959note} explora el espacio de búsqueda manteniendo una partición de los vértices en dos conjuntos disjuntos: el conjunto $S$ de vértices cuya distancia mínima desde el origen $s$ ya ha sido permanentemente determinada, y el conjunto $Q = V \setminus S$ de vértices tentativos gestionados en una cola de prioridad. En cada paso iterativo, el algoritmo extrae el vértice $u \in Q$ con la menor distancia tentativa acumulada $d[u]$, lo incorpora a $S$, y relaja cada arista saliente $(u, v) \in E$:
\begin{equation}
    \text{si } d[u] + w(u, v) < d[v] \implies d[v] \leftarrow d[u] + w(u, v), \quad \pi[v] \leftarrow u
\end{equation}
La actualización del valor $d[v]$ en la cola de prioridad exige ejecutar la operación \texttt{Decrease-Key}$(Q, v, d[v])$.

\subsubsection{Estructuras de Colas de Prioridad con Decrease-Key}
La cota asintótica de Dijkstra está dominada por las tres primitivas de la cola:
\begin{itemize}
    \item \textbf{Min-Heap Binario Indexado:} Árbol binario casi completo representado en un arreglo contiguo en memoria, complementado por una tabla hash inversa de posiciones $\text{pos}: V \to \mathbb{N}$ que ubica cualquier vértice en $O(1)$. \texttt{Insert} toma $O(\log |V|)$, \texttt{Extract-Min} toma $O(\log |V|)$ y \texttt{Decrease-Key} toma $O(\log |V|)$.
    \item \textbf{4-ary Heap Contiguo ($d=4$):} Árbol $d$-ario con factor de ramificación 4. Cada nodo padre en el índice $i$ tiene sus 4 hijos en los índices $4i+1, 4i+2, 4i+3, 4i+4$. Esto reduce la profundidad del árbol a la mitad ($\log_4 |V| = 0.5 \log_2 |V|$), requiriendo menos niveles de flotado en \texttt{Decrease-Key} y aprovechando las líneas de caché L1/L2 de 64 bytes.
    \item \textbf{Fibonacci Heap:} Colección de árboles con orden de montículo estructurados mediante listas circulares doblemente enlazadas con punteros \texttt{parent}, \texttt{child}, \texttt{left}, \texttt{right}, \texttt{degree} y \texttt{marked} \parencite{fredman1987fibonacci}. Utiliza inserción y corte perezoso (\textit{lazy}), permitiendo que \texttt{Insert} y \texttt{Decrease-Key} operen en $O(1)$ amortizado, y \texttt{Extract-Min} en $O(\log |V|)$ amortizado mediante consolidación de árboles de igual grado.
\end{itemize}

\subsubsection{Extensión Algorítmica: Búsqueda Heurística $A^*$ y Poda Espacial}
Una limitación intrínseca de Dijkstra en consultas punto a punto ($s \to t$) es que expande nodos en forma de onda radial concéntrica en todas direcciones. El Algoritmo $A^*$ \parencite{hart1968formal} focaliza la búsqueda mediante una función de evaluación compuesta:
\begin{equation}
    f(n) = g(n) + h(n)
\end{equation}
donde $g(n)$ es el costo acumulado desde el origen $s$, y $h(n)$ es una heurística admisible y consistente que estima el costo residual hacia el destino $t$. Para garantizar optimalidad, la heurística debe satisfacer:
\begin{enumerate}
    \item \textbf{Admisibilidad:} $0 \le h(n) \le h^*(n)$ para todo $n \in V$, donde $h^*(n)$ es el costo óptimo real.
    \item \textbf{Consistencia (Monotonicidad):} $h(u) \le w(u, v) + h(v)$ para toda arista dirigida $(u, v) \in E$.
\end{enumerate}
En este proyecto, para la red vial cusqueña con pesos de tiempo, la heurística euclidiana se escala mediante el factor admisible:
\begin{equation}
    \lambda = \min_{(u,v) \in E} \frac{w(u,v)}{\|x_u - x_v\|}, \quad h(u) = \lambda \cdot \|x_u - x_t\|
\end{equation}
asegurando consistencia matemática rigurosa y podando expansiones sin alterar el costo mínimo óptimo.

\subsubsection{Metodología de Pruebas Unitarias: El Patrón AAA}
La totalidad de la suite de pruebas del proyecto (278 pruebas con Pytest) se implementa bajo el patrón \textbf{Arrange - Act - Assert (AAA)}:
\begin{itemize}
    \item \textbf{Arrange:} Inicialización de precondiciones, generación determinista de instancias con semilla fija y selección de la cola a evaluar.
    \item \textbf{Act:} Invocación del solver (\texttt{dijkstra\_solver} o \texttt{a\_star\_solver}) y captura de vectores de distancias y contadores de operaciones.
    \item \textbf{Assert:} Verificación estricta contra oráculos independientes (Bellman-Ford, NetworkX) y evaluación del certificado de optimalidad.
\end{itemize}

\subsection{Fuentes Estudiadas}
La selección de literatura se condujo con criterios de rigor académico, priorizando artículos seminales arbitrados por pares, manuales algorítmicos estándar de la Association for Computing Machinery (ACM) e IEEE, e investigaciones sobre análisis de rendimiento bajo jerarquía de memoria (Tabla \ref{tab:criterios_fuentes}).

\begin{table}[H]
    \centering
    \fontsize{9.5}{11.5}\selectfont
    \renewcommand{\arraystretch}{1.15}
    \caption{Selección razonada de fuentes bibliográficas clave y criterios de pertinencia}
    \label{tab:criterios_fuentes}
    \begin{tabularx}{\textwidth}{@{}>{\raggedright\arraybackslash}p{3.6cm} >{\raggedright\arraybackslash}p{3.8cm} >{\raggedright\arraybackslash}X@{}}
        \toprule
        \textbf{Fuente Académica} & \textbf{Tipología y Dominio} & \textbf{Criterio de Pertinencia e Impacto en el Proyecto} \\
        \midrule
        \textcite{dijkstra1959note} & Artículo Seminal (Numerische Mathematik) & Fundamento teórico del algoritmo SSSP con pesos no negativos y principio de relajación voraz. \\
        \textcite{fredman1987fibonacci} & Artículo Teórico (Journal of the ACM) & Especificación formal del Fibonacci Heap con tiempos amortizados $O(1)$ para \texttt{Decrease-Key}. \\
        \textcite{tarjan1983data} & Monografía Clásica (SIAM) & Análisis riguroso del método del potencial para el cálculo de costos amortizados en estructuras de datos. \\
        \textcite{cherkassky1996shortest} & Estudio Empírico (Journal of Algorithms) & Evaluación experimental exhaustiva de algoritmos de caminos mínimos; evidencia de ineficiencias del Fibonacci Heap en grafos dispersos. \\
        \textcite{lamarca1996cache} & Arquitectura y Algoritmos (ACM SODA) & Análisis del impacto de las líneas de caché L1/L2 y fallos de traducción TLB en montículos binarios y $d$-arios. \\
        \textcite{hart1968formal} & Artículo Seminal (IEEE SSC) & Formulación de $A^*$ y condiciones de admisibilidad y consistencia para poda de grafos. \\
        \textcite{boeing2017osmnx} & Sistemas Geoespaciales (Computers, Env. and Urban Systems) & Metodología de ingestión, validación y corrección topológica de grafos urbanos a partir de OpenStreetMap. \\
        \textcite{cormen2022introduction} & Libro de Texto Avanzado (MIT Press) & Demostración formal de invariantes de bucle, cotas de recurrencia y pseudocódigo estructurado. \\
        \bottomrule
    \end{tabularx}
\end{table}

\subsection{Conocimientos Nuevos y Brechas Técnicas}
Al inicio del proyecto, el equipo diagnosticó cuatro brechas conceptuales y técnicas que trascendían los contenidos estándar de cursos formativos previos:
\begin{enumerate}
    \item \textbf{Brecha 1: Manipulación de Estructuras Arbóreas con Punteros Múltiples.} Dominio de la representación en memoria de listas circulares doblemente enlazadas con punteros padre, hijo, izquierdo y derecho para el Fibonacci Heap, garantizando concatenaciones en $O(1)$ sin pérdida de referencias.
    \item \textbf{Brecha 2: Operación Decrease-Key en Tiempo Constante Amortizado $O(1)$.} Comprensión analítica del mecanismo de marcado de nodos (\textit{node marking}) y cortes en cascada (\textit{cascading cuts}) propuesto por Tarjan, asegurando que el grado de cualquier nodo en un heap de $n$ elementos no exceda $\lfloor \log_\phi n \rfloor$.
    \item \textbf{Brecha 3: Arquitectura de Computadores, Jerarquía de Memoria y Líneas de Caché L1/L2.} Comprensión del fenómeno por el cual una estructura asintóticamente superior como el Fibonacci Heap es superada en hardware contemporáneo por estructuras con peor cota teórica ($d$-ary heap) debido a la contigüidad espacial en los bloques de caché de 64 bytes.
    \item \textbf{Brecha 4: Modelado de Redes Viales Reales con OpenStreetMap y Grafos Geoespaciales.} Extracción de mallas viales urbanas reales mediante OSMnx, conversión de primitivas geoespaciales WGS84 a grafos dirigidos ponderados y manejo de tramos viales paralelos y sentidos únicos.
\end{enumerate}

\subsection{Estrategia y Evidencias de Aprendizaje Autónomo}
Para cerrar las brechas identificadas sin supervisión externa continua, el equipo adoptó una estrategia estructurada en cuatro hitos:
\begin{enumerate}
    \item \textbf{Hito 1 (Investigación y Prototipo Temprano):} Lectura directa de \textcite{cormen2022introduction} y \textcite{fredman1987fibonacci}. Implementación de un prototipo temprano del Min-Heap binario indexado con tabla de posiciones $O(1)$ para validar la mecánica de \texttt{decrease\_key}.
    \item \textbf{Hito 2 (Extensión d-aria y Análisis de Caché):} Estudio de la investigación de \textcite{lamarca1996cache} e implementación del montículo 4-ario contiguo, verificando que la reducción de niveles en el árbol compensa el aumento de comparaciones por nivel.
    \item \textbf{Hito 3 (Fibonacci Heap y Detección de Defectos):} Construcción manual del montículo de Fibonacci con punteros dobles circulares. Durante las pruebas de estrés, el equipo detectó de forma autónoma que un arreglo de consolidación con tamaño estático $\log_2(n)+7$ causaba pérdida de árboles; se corrigió analíticamente redimensionándolo dinámicamente con la cota de la razón áurea $\lceil \log_\phi(n) \rceil + 2$.
    \item \textbf{Hito 4 (Validación Rigurosa):} Implementación del certificador de optimalidad de Bellman $O(|V|+|E|)$ y de la suite automatizada de 278 pruebas con Pytest, asegurando que las tres estructuras produzcan distancias matemáticamente idénticas en todas las instancias.
\end{enumerate}
Las evidencias de este aprendizaje autónomo se encuentran registradas cronológicamente en los commits del repositorio de GitHub y en la suite de pruebas unitarias.
